% 由既有正文整理；包括标点符号不超过400字。
本课题面向新疆复杂地形下气候模式粗网格与区域气象信息需求之间的尺度差距，研究CIESM近地面气温、风场、短波辐射和降水的多变量时空超分辨率与尺度一致性校正。针对细网格参考场多由插值生成、可学信号未经检验的问题，先逐变量检验候选目标场的真实信息增量并据此选源，再构造粗细配对样本，结合地形条件与多尺度时间表征，比较联合与等信息独立模型。将聚合一致性、非负及变量定义条件与经验性地形约束分开处理，形成可逐项检查的TerraSR-XJ约束投影模块。评价采用连续年份划分与地形留出，并以不参与训练的国家级台站日值观测作独立验证，判定阈值在查看测试集前写定。预期交付算法模块、新疆0.1度逐日四变量试验数据产品与可复现评估流程，实现成果与CIESM集成并实际应用，在本领域重要学术期刊发表或获接收论文不少于1篇。
